\documentclass[10pt,letterpaper,oneside]{article}
\usepackage[T1]{fontenc}
\usepackage[utf8]{inputenc}
\usepackage{lmodern}
\usepackage{geometry}
\geometry{margin=0.84in,headheight=21pt,footskip=28pt}
\usepackage{microtype}
\usepackage{xcolor}
\definecolor{ink}{HTML}{15253D}
\definecolor{accent}{HTML}{086A84}
\definecolor{pale}{HTML}{EEF5F8}
\definecolor{muted}{HTML}{4E6474}
\definecolor{grayline}{HTML}{D7E2E8}
\usepackage{graphicx}
\usepackage{fancyhdr}
\usepackage{titlesec}
\usepackage{enumitem}
\usepackage{array,tabularx,longtable,booktabs}
\usepackage[most]{tcolorbox}
\usepackage{xurl}
\usepackage{hyperref}
\usepackage{listings}
\usepackage{amsmath,amssymb}
\usepackage{tikz}
\usetikzlibrary{arrows.meta,positioning,shapes.geometric,calc}
\hypersetup{colorlinks=true,linkcolor=accent,urlcolor=accent,pdftitle={Vulnerability Assessment, Forensics and Incident Response},pdfauthor={Independent technical engineering study companion}}
\renewcommand{\familydefault}{\sfdefault}
\linespread{1.07}
\setlength{\parindent}{0pt}
\setlength{\parskip}{5.5pt}
\emergencystretch=2em
\setlist[itemize]{leftmargin=1.5em,itemsep=2.5pt,topsep=3pt}
\setlist[enumerate]{leftmargin=1.7em,itemsep=3pt,topsep=3pt}
\titleformat{\section}{\Large\bfseries\color{ink}}{\thesection}{0.65em}{}[\vspace{-2pt}{\color{accent}\titlerule}]
\titleformat{\subsection}{\large\bfseries\color{accent}}{\thesubsection}{0.6em}{}
\titleformat{\subsubsection}{\bfseries\color{ink}}{\thesubsubsection}{0.6em}{}
\titlespacing*{\section}{0pt}{16pt}{9pt}
\titlespacing*{\subsection}{0pt}{11pt}{6pt}
\pagestyle{fancy}\fancyhf{}\fancyhead[L]{\footnotesize\color{muted} COMPUTER \& INFORMATION SECURITY ENGINEERING}\fancyhead[R]{\footnotesize\color{accent} SUITE / 6}\fancyfoot[L]{\footnotesize\color{muted} Independent study companion \quad|\quad 2026}\fancyfoot[R]{\footnotesize\color{muted}\thepage}\renewcommand{\headrulewidth}{0.4pt}\renewcommand{\headrule}{\hbox to\headwidth{\color{grayline}\leaders\hrule height \headrulewidth\hfill}}
\newtcolorbox{sourcebox}{enhanced,breakable,boxrule=0pt,leftrule=3pt,colframe=accent,colback=pale,arc=0pt,left=9pt,right=7pt,top=7pt,bottom=7pt}
\newtcolorbox{warningbox}{enhanced,breakable,boxrule=0.5pt,colframe=grayline,colback=white,arc=1pt,left=9pt,right=9pt,top=7pt,bottom=7pt}
\newcommand{\handbookref}[1]{\textsf{\color{accent}\textbf{[H:#1]}}}
\newcommand{\modernext}{\textsf{\color{muted}\textbf{[E]}}}
\newcommand{\designed}{\textsf{\color{accent}\textbf{[D]}}}
\newcommand{\deliverable}[1]{\textbf{Deliverable:} #1}
\lstdefinestyle{lab}{basicstyle=\ttfamily\scriptsize,keywordstyle=\color{accent}\bfseries,commentstyle=\color{muted},numbers=left,numberstyle=\tiny\color{muted},numbersep=7pt,breaklines=true,showstringspaces=false,frame=single,rulecolor=\color{grayline},backgroundcolor=\color{pale},xleftmargin=1em,framexleftmargin=0.6em,tabsize=2,columns=fullflexible}
\lstset{style=lab}
\renewcommand{\arraystretch}{1.21}
\newcommand{\makereportcover}[4]{%
\thispagestyle{empty}\vspace*{1.15in}
{\small\bfseries\color{accent} COMPUTER AND INFORMATION SECURITY ENGINEERING}\par\vspace{13pt}
{\Huge\bfseries\color{ink} #1\par}\vspace{15pt}
{\Large\color{muted} #2\par}\vspace{18pt}
\color{accent}\rule{0.94\textwidth}{2.1pt}\par\vspace{20pt}
{\normalsize\bfseries Document #3 of 8 \quad\textbullet\quad 16 October 2026\par}
\vspace{0.35in}
\begin{sourcebox}\textbf{Orientation.} #4\end{sourcebox}
\vfill
{\small\color{muted} Independent, original engineering study and practice materials grounded in\par
\emph{Computer and Information Security Handbook}, 4th ed., edited by John R. Vacca,\par
Elsevier/Morgan Kaufmann, 2025, Volumes I--II. Not a reproduction or official standard.}
\newpage\setcounter{page}{1}\tableofcontents\newpage
}
\begin{document}
\makereportcover{Vulnerability Assessment, Forensics and Incident Response}{Risk-based triage, application assessment, evidence, response and assurance}{6}{Integrate the handbook's vulnerability assessment and forensic investigation foundations with modern secure-SDLC analysis and defensible incident-response evidence.}

\section{Assessment scope, authority and intended outcomes}
The handbook distinguishes vulnerability assessment from penetration testing, highlights scanner validation and covers firmware, cyber forensics, incident investigation, e-discovery, network traces, metadata and disk imaging. \handbookref{31,32,41--46,77,81} Modern scanning and software security practices below are extensions, not retroactively attributed to the text.
\begin{sourcebox}\textbf{Assessment charter} \designed{}: sponsor and authorization; systems in scope; allowed techniques and rate limits; prohibited actions; data handling; points of contact; emergency stop; evidence retention; output classification; remediation ownership; date and version. Never interpret an audit request as permission to exploit production systems.\end{sourcebox}

\section{Vulnerability assessment workflow}
\begin{enumerate}
\item Inventory assets, owners, technologies, versions and business criticality.
\item Establish authenticated/unauthenticated scan scope and exclusions.
\item Verify credentials, agent health, coverage and scan safety on representative assets.
\item Capture findings and raw evidence with tool versions and detection timestamps.
\item De-duplicate by affected asset/component/root cause without erasing instances.
\item Validate reachability, exploitability evidence, compensating controls and false-positive possibility.
\item Prioritize and assign fixes, mitigations or documented exceptions.
\item Retest the \emph{deployed} fix, verify no regressions and record closure evidence.
\end{enumerate}
\handbookref{31,33,35,81} \designed{}
\subsection{Signal taxonomy and prioritization}\modernext{}
\par\medskip\noindent\begin{tabularx}{\textwidth}{@{}p{0.19\textwidth}X@{}}\toprule
Signal & Interpretation / limits \\ \midrule
CVE / GHSA & Identifier of a disclosed vulnerability or advisory; not a measured risk.\\
CWE & Weakness category linking root cause to remediation patterns.\\
CVSS & A severity scoring system; not a substitute for asset context.\\
CISA KEV & Evidence a CVE is known to have been exploited in the wild; prioritization input.\\
EPSS & Statistical likelihood estimate of exploitation in a time horizon; not certainty.\\
Exposure & Path, privilege and controls affecting whether attacker reaches vulnerable functionality.\\
Observed exploit & Confirmed activity or defensible demonstration; stronger than theoretical path.\\
Business impact & Material consequence to customer, operational, safety or legal obligations.\\
\bottomrule\end{tabularx}\par\medskip
CISA explicitly presents KEV as an input to vulnerability prioritization. Do not equate absence from KEV with safety. [C1]

\section{Application security assessment methods}\handbookref{10,31,77}\modernext{}
\subsection{Data flow, sinks and trust assumptions}
Modern static application security testing may use syntax pattern matching, interprocedural dataflow, taint tracking and semantic models. A taint-flow result requires analysis of \emph{source, transformation, sanitization, sink and path feasibility}. A path that appears in source code may not be reachable or exploitable in deployment.
\par\medskip\noindent\begin{tabularx}{\textwidth}{@{}p{0.2\textwidth}X X@{}}\toprule
Technique & Typical evidence & Major blind spot \\ \midrule
SAST / CodeQL & Code path, rule and source-to-sink trace & Runtime configuration and missing models\\
SCA / dependency scan & Package/version to advisory match & Contextual reachability and build provenance\\
Secret scanning & Credential-like text and provider pattern & Validity, privilege and exposure context\\
DAST & HTTP request/response showing misbehavior & Inaccessible paths and environment mismatch\\
IAST & Instrumented runtime execution finding & Coverage depends on exercised flows\\
Manual review & Business-logic violation and design context & Reviewer/time constraints\\
\bottomrule\end{tabularx}\par\medskip
\subsection{A reproducible finding record}\designed{}
Record: ID; application/repo; branch and commit; deployment environment; affected endpoint, file or dependency; CWE/CVE; source, sink and taint path if relevant; impact scenario; preconditions; observed evidence; reproduction in authorized scope; confidence; exposure; suggested fix; owner; agreed deadline; retest and final disposition.

\section{Modern repository security controls}\modernext{}
NIST SSDF v1.1 recommends secure software development practices that can be integrated into diverse SDLC processes. Its publication is guidance, not an automatic legal requirement. [N2]
\begin{itemize}
\item Run analysis on pull requests and relevant protected branches; retain tool version and commit SHA.
\item For blocking checks, confirm analysis ran on the \emph{candidate merge commit} or expected branch/commit, not only an unrelated historical scan.
\item Review detected secrets for validity and permissions; rotate true exposures before declaring remediation complete.
\item Generate and retain machine-readable findings such as SARIF 2.1.0 where tool ecosystem supports it.
\item Link the deployed remediation to the original alert and verify it in a suitable runtime.
\item Use documented triage reasons; do not mechanically label unverified findings ``false positive''.
\end{itemize}
GitHub Code Scanning supports alert retrieval and SARIF submission via its REST API; access and availability depend on repository configuration and permissions. [G1]

\section{Forensics: preserve integrity and provenance}
\handbookref{41--46} Evidence acquisition is a controlled process distinct from ordinary troubleshooting. Before using a tool, identify legal authority, data ownership, scope restrictions, volatility, retention and confidentiality.
\subsection{Evidence handling record}\designed{}
\par\medskip\noindent\begin{tabularx}{\textwidth}{@{}p{0.25\textwidth}X@{}}\toprule
Field & Example / expectation \\ \midrule
Evidence ID & CASE-014/EV-003\\
Origin & Hostname/service/volume/log source and system owner\\
Acquired by & Investigator and authorizing incident lead\\
Collection time & UTC timestamp, time source, recorded clock skew\\
Method & Acquisition tool version, commands and read-only precautions\\
Integrity & SHA-256 digest on acquisition and on access/transfer\\
Custody & Each transfer: person, time, purpose, storage and signature\\
Restrictions & PII, attorney-client, regulatory and retention rules\\
\bottomrule\end{tabularx}\par\medskip
\subsection{Scope of forensic methods}
\begin{itemize}
\item \textbf{Disk imaging:} preserve original media state; distinguish HDD from SSD/trim behavior; use compatible forensic acquisition methods.
\item \textbf{Memory:} collect volatile artifacts where authorized and appropriate, noting acquisition side effects.
\item \textbf{Network:} capture payload/flow/metadata under approved legal and privacy scope.
\item \textbf{Email and documents:} examine metadata and trace headers carefully; claims of origin require corroboration.
\item \textbf{Cloud:} snapshot relevant resource and control-plane records; document provider time semantics and shared responsibility.
\end{itemize}

\section{Incident response lifecycle and coordination}\modernext{}
NIST SP 800-61 Rev.~3 (2025) places incident response in enterprise cyber risk management aligned to CSF 2.0. The handbook's forensic and incident chapters provide investigative foundations. \handbookref{42,44} [N3]
\begin{enumerate}
\item Confirm and classify incident, with known facts versus hypotheses clearly separated.
\item Declare incident lead, technical teams, legal/privacy coordination and communications authority.
\item Preserve critical evidence before actions that would destroy volatile state, where feasible and safe.
\item Contain with least-disruptive effective means; validate that containment actually worked.
\item Recover through trusted sources; rotate secrets and rebuild assets when integrity is uncertain.
\item Confirm business service, data integrity and detection health; communicate residual exposure.
\item Hold blameless technical review and track systemic fixes to completion.
\end{enumerate}
\subsection{Example timeline analysis}\designed{}
\par\medskip\noindent\begin{tabularx}{\textwidth}{@{}p{0.17\textwidth}X p{0.26\textwidth}@{}}\toprule
UTC time & Observation & Confidence / corroboration \\ \midrule
09:02 & Unusual privileged token used & Identity log; high confidence\\
09:04 & Storage list API called & Cloud audit; high confidence\\
09:06 & Large outbound transfer observed & Flow log; medium confidence; payload unknown\\
09:11 & Token revoked & IAM record; high confidence\\
09:18 & Requests denied after revocation & Gateway tests; high confidence\\
\bottomrule\end{tabularx}\par\medskip
Do not claim confirmed exfiltration based only on a traffic-volume anomaly; validate what actually left, if evidence permits.

\section{Closure, remediation and audit trail}
\begin{warningbox}\textbf{Closure criteria} \designed{}: root cause identified (or uncertainty acknowledged); affected assets and identities scoped; fix or compensating control deployed; regression test passed; incident data handled per policy; evidence archived with custody; risk owner accepts remaining exposure; detection and prevention improvements assigned.\end{warningbox}

\section{Case exercises and expected outputs}
\textbf{Case A:} vulnerable library exists but the affected feature is disabled. Supply a triage rationale, proof and a plan for future exposure monitoring. \textbf{Case B:} secret scanning reports a public client identifier. Differentiate identifier from privileged secret, prove provider semantics and document disposal/triage. \textbf{Case C:} endpoint telemetry indicates malicious execution; design the first-hour evidence and containment workflow with confidentiality safeguards. \designed{}
\begin{sourcebox}\textbf{Assessment rubric (100):} authorized scope (15), reproducible evidence (30), risk reasoning (20), forensics and custody (20), verified remediation (15).\end{sourcebox}

\section*{Handbook chapter alignment}\addcontentsline{toc}{section}{Handbook chapter alignment}
These chapters ground the domain scope. The document offers original synthesis and technical procedures, not a chapter-by-chapter reproduction.
\par\medskip\noindent\begin{longtable}{@{}p{0.085\textwidth}p{0.68\textwidth}p{0.16\textwidth}@{}}\toprule Chapter & Handbook topic & Combined PDF p. \\ \midrule\endhead10 & Securing Web Applications, Services, and Servers & 211 \\
28 & The Enemy (The Intruder's Genesis) & 545 \\
29 & Social Engineering Deceptions and Defenses & 557 \\
30 & Hacking in an Ethical Way & 567 \\
31 & What Is Vulnerability Assessment? & 573 \\
32 & Introduction to Vulnerabilities Below the Operating System & 585 \\
33 & Security Metrics: An Introduction and Literature Review & 595 \\
35 & Risk Management & 619 \\
41 & Cyber Forensics & 687 \\
42 & Cyber Forensics and Incidence Response & 717 \\
43 & Securing e-Discovery & 743 \\
44 & Network Forensics & 769 \\
45 & Microsoft Office and Metadata Forensics: A Deeper Dive & 785 \\
46 & Hard Drive Imaging & 795 \\
56 & Detection of Conflicts in Security Policies & 1003 \\
77 & Penetration Testing & 1337 \\
81 & Assessments and Audits & 1369 \\
87 & Cyber Warfare & 1431 \\
88 & Cyber Attack Process & 1467 \\
97 & Data Loss Protection & 1629 \\
101 & Use of Artificial Intelligence (AI) in Cybersecurity & 1703 \\
102 & AI Abuse: Investigating the Threat Landscape & 1711 \\
\bottomrule\end{longtable}\par\medskip

\section*{Source and reference register}
\addcontentsline{toc}{section}{Source and reference register}
\begin{description}[leftmargin=1.6em,style=nextline]
\item[H] John R. Vacca (ed.), \emph{Computer and Information Security Handbook}, 4th ed., Volumes I--II, Elsevier/Morgan Kaufmann, 2025. Chapter identifiers in this document point to this book.
\item[N1] NIST, \emph{Cybersecurity Framework 2.0} (2024), \url{https://www.nist.gov/cyberframework}.
\item[N2] NIST SP 800-218, \emph{Secure Software Development Framework v1.1} (2022), \url{https://csrc.nist.gov/pubs/sp/800/218/final}. Later drafts are not treated as final requirements.
\item[N3] NIST SP 800-61 Rev. 3, \emph{Incident Response Recommendations and Considerations for Cybersecurity Risk Management} (2025), \url{https://csrc.nist.gov/pubs/sp/800/61/r3/final}.
\item[N4] NIST SP 800-63-4, \emph{Digital Identity Guidelines} (2025), \url{https://csrc.nist.gov/pubs/sp/800/63/4/final}.
\item[N5] NIST, FIPS 203/204/205 (2024), \url{https://csrc.nist.gov/Projects/Post-Quantum-Cryptography/publications}.
\item[O1] OWASP, \emph{API Security Top 10} (2023), \url{https://api-security.owasp.org/editions/2023/en/0x11-t10/}.
\item[C1] CISA, \emph{Known Exploited Vulnerabilities Catalog}, \url{https://www.cisa.gov/known-exploited-vulnerabilities-catalog}.
\item[M1] MITRE ATT\&CK Enterprise, \url{https://attack.mitre.org/matrices/enterprise/}.
\item[G1] GitHub, \emph{REST API endpoints for code scanning}, \url{https://docs.github.com/en/rest/code-scanning}.
\end{description}
\begin{warningbox}\textbf{Interpretation rule.} \handbookref{number} identifies a conceptual anchor in the supplied handbook, not a verbatim quotation or an official requirement. \modernext{} identifies a contemporary extension informed by external references or current engineering practice. \designed{} identifies an independently developed example, template, exercise, metric, or operational procedure. Requirements are binding only if separately adopted through applicable law, contract, policy, or the cited standards; this suite does not assert otherwise.\end{warningbox}

\end{document}
